\documentclass[12pt]{article}
\usepackage[utf8]{inputenc}
\usepackage[T1]{fontenc}
\usepackage{hyperref}
\title{Torau Language Resources for K-12 Curriculum Development}
\author{Compiled by Claude Sonnet 5 (claude-sonnet-5) for Andr\'as Kornai}
\date{}
\begin{document}
\maketitle
\section{Introduction}
Torau (ISO 639-3 ttu; Glottocode tora1259) is an Austronesian language of the Mono-Uruavan subgroup of Northwest Solomonic. It is spoken on the east coast of Bougainville, Papua New Guinea, in Rorovana, Vito and Tarara. Ethnologue reported 605 speakers in 1963, and one source estimates about 1,200 today. The language is highly endangered. The only substantial resources found are linguistic documentation.
\section{Primary Educational Resources}
No textbooks, primers or school programmes were found in search.
\section{Linguistic and Reference Materials}
\begin{itemize}
\item Torau overview, Surrey Morphology Group (Northwest Solomonic project): \url{https://www.smg.surrey.ac.uk/languages/northwest-solomonic/torau/overview/}
\item Torau dictionary (work in progress, Bill Palmer): \url{https://www.smg.surrey.ac.uk/languages/northwest-solomonic/torau/torau-dictionary/}
\item Wikipedia, Torau language: \url{https://en.wikipedia.org/wiki/Torau_language}
\item PARADISEC item (Capell papers, Torau and Uruava notes): \url{https://catalog.paradisec.org.au/repository/AC2/NSBURU102}
\end{itemize}
Before the current documentation work, the only grammatical description was a three-page set of notes by Peter Rausch (1912).
\section{Religious and Community Resources}
None found in search.
\section{Digital and Audio Resources}
Search summaries say a text corpus with audio and interlinear transcripts is posted on the Surrey Morphology Group site. I did not retrieve a direct link to it, so see the overview page above.
\section{Curriculum Development Context}
Torau is not known to be taught in schools. The dictionary and corpus could supply vocabulary and texts, but they were made for linguistic documentation and not for classroom use. Any curriculum would need to be built with the community.
\end{document}